\section{Protocolo experimental y correcciones}

\subsection{Condiciones de comparación}

Se usan los mismos 12,753 ejemplos reales de entrenamiento y 3,712 de validación,
104 clases, entrada $112\times112$, lote de 64 y semilla 0. Se permiten hasta 150
épocas. Se guarda el checkpoint con el mayor F1 macro de validación; por ello,
la exactitud final reportada corresponde a ese mismo checkpoint y no al máximo
independiente de exactitud. La referencia totalmente conectada y las dos CNN
comparadas por profundidad comparten optimizador, aumentos, dropout y paciencia
de cinco evaluaciones sin mejora. La intervención combinada de regularización
permite una paciencia mayor y cambia el optimizador y su tasa; no es una ablación
aislada de cada técnica.

Las curvas contienen métricas de entrenamiento con aumentos y dropout activos,
mientras que la validación se calcula sin ellos. Por esta razón, una exactitud
de validación mayor que la de entrenamiento no implica por sí sola fuga de datos.
Para diagnosticar sobreajuste se examinan conjuntamente la tendencia de la
pérdida, la exactitud y F1 macro, sin usar únicamente la distancia entre dos curvas.

La comparación se limita a una semilla y utiliza validación para seleccionar
épocas e hiperparámetros. No proporciona intervalos de confianza ni una medida
final independiente de generalización. El conjunto de prueba de Kaggle carece
de etiquetas locales; no se presentan los resultados de validación como
exactitud de prueba ni como puntaje de la competencia.

La semilla y las configuraciones se guardan para repetir el procedimiento, pero
las operaciones de GPU y la canalización paralela no garantizan resultados
idénticos bit a bit. Las métricas históricas y las nuevas ejecuciones se identifican
por separado en sus respectivas carpetas.

\subsection{Trazabilidad y ajustes del documento}

Las ejecuciones nuevas se guardan bajo \texttt{reports/avance4\_2026-10-21}, con
configuraciones, históricos por época, checkpoints y predicciones por imagen.
Las figuras se exportan tanto a PNG para presentación como a PDF para el reporte.
Un manifiesto registra la versión del código y sus hashes; los experimentos
anteriores se conservan y no se mezclan con las métricas nuevas.

Respecto del documento anterior, las secciones 1 a 6 ahora explicitan los
objetivos medibles, la resolución realmente usada y el lote de 64; representan
las arquitecturas y la ubicación exacta de las dos capas nuevas; distinguen la
referencia totalmente conectada de la CNN original; separan destilación y
clasificación; y aclaran el criterio de selección y el alcance de la validación.
Se atiende la observación compartida sobre las representaciones de los modelos.
No se dispone del texto completo de la retroalimentación del Avance 3, por lo que
no se afirma haber atendido observaciones adicionales no documentadas.
